\documentclass[a4paper,11pt]{article}

\usepackage[a4paper,margin=2cm]{geometry}
\usepackage[T1]{fontenc}
\usepackage[utf8]{inputenc}
\usepackage[ngerman,english]{babel}
\usepackage{lmodern}
\usepackage{microtype}
\usepackage{xcolor}
\usepackage{colortbl}
\usepackage{array}
\usepackage{parskip}
\usepackage{enumitem}
\usepackage[most]{tcolorbox}
\usepackage{titlesec}
\usepackage[hidelinks]{hyperref}

\definecolor{HeaderBlueDark}{RGB}{70,110,170}
\definecolor{RowBlueLight}{RGB}{235,243,255}
\definecolor{RowBlueDark}{RGB}{220,233,250}
\definecolor{SoftGray}{RGB}{90,90,90}

\setlength{\parindent}{0pt}
\setlist[itemize]{leftmargin=1.4em,itemsep=0.35em,topsep=0.35em}
\linespread{1.06}
\renewcommand{\arraystretch}{1.35}

\titleformat{\section}[block]
{\Large\bfseries\color{HeaderBlueDark}}
{}{0pt}{}
\titlespacing{\section}{0pt}{1.3em}{0.8em}

\titleformat{\subsection}[block]
{\large\bfseries\color{HeaderBlueDark}}
{}{0pt}{}
\titlespacing{\subsection}{0pt}{1.0em}{0.6em}

\newtcolorbox{exbox}{enhanced,breakable,colback=RowBlueLight,colframe=RowBlueDark,boxrule=0.4pt,arc=1.5mm,left=1.2mm,right=1.2mm,top=1mm,bottom=1mm,title={Beispiele \textnormal{(Examples)}},fonttitle=\bfseries,colbacktitle=RowBlueDark,coltitle=black}

\NewDocumentEnvironment{grammarentry}{mm+b}{%
    \begin{tcolorbox}[enhanced,breakable,colback=white,colframe=HeaderBlueDark,boxrule=0.6pt,arc=1.5mm,left=1.5mm,right=1.5mm,top=1mm,bottom=1mm,colbacktitle=HeaderBlueDark,coltitle=white,fonttitle=\bfseries,title={#1\hfill\normalfont\footnotesize #2}]
    #3
    \end{tcolorbox}
}{}

\newcommand{\GermanText}[1]{\textbf{Deutsch: }#1\par}
\newcommand{\EnglishText}[1]{{\small\color{SoftGray}\textbf{English: }#1\par}}
\newcommand{\ExampleItem}[2]{\item \textbf{#1}\\{\small\color{SoftGray}#2}}

\newcommand{\HeaderRow}{\rowcolor{HeaderBlueDark}\color{white}\bfseries Fall & \color{white}\bfseries Maskulin & \color{white}\bfseries Feminin & \color{white}\bfseries Neutrum & \color{white}\bfseries Plural \\}
\newcommand{\DeclTable}[4]{%
\begin{center}
\begin{tabular}{|>{\centering\arraybackslash}m{2.2cm}|>{\centering\arraybackslash}m{2.5cm}|>{\centering\arraybackslash}m{2.5cm}|>{\centering\arraybackslash}m{2.5cm}|>{\centering\arraybackslash}m{2.5cm}|}
\hline
\HeaderRow
\hline
\rowcolor{RowBlueLight}Nominativ & #1 \\
\hline
\rowcolor{RowBlueDark}Akkusativ & #2 \\
\hline
\rowcolor{RowBlueLight}Dativ & #3 \\
\hline
\rowcolor{RowBlueDark}Genitiv & #4 \\
\hline
\end{tabular}
\end{center}}

\title{Demonstrativpronomen und verwandte Formen}
\date{}

\begin{document}
\maketitle
\tableofcontents
\newpage

\section{Überblick}
\begin{grammarentry}{Grammatische Rolle}{Pronomen und Artikelwörter}
\GermanText{Wörter wie \textbf{dieser}, \textbf{jener}, \textbf{derselbe} und \textbf{derjenige} können als Demonstrativpronomen verwendet werden. Sie weisen besonders auf eine bestimmte Person, Sache oder Gruppe hin.}
\EnglishText{Words such as \textbf{dieser}, \textbf{jener}, \textbf{derselbe}, and \textbf{derjenige} can be used as demonstrative pronouns. They point specifically to a particular person, thing, or group.}
\GermanText{Stehen sie vor einem Nomen, werden sie oft als Demonstrativartikel oder demonstrative Artikelwörter bezeichnet.}
\EnglishText{When they stand before a noun, they are often called demonstrative determiners or demonstrative article words.}
\end{grammarentry}

\section{Dieser und jener}
\begin{grammarentry}{dieser / jener}{Demonstrativpronomen, deklinierbar}
\GermanText{\textbf{dieser} verweist meist auf etwas Näheres; \textbf{jener} verweist meist auf etwas Weiteres oder zeitlich Entferntes. \textbf{jener} wirkt häufig schriftsprachlich oder gehoben.}
\EnglishText{\textbf{dieser} usually refers to something nearer; \textbf{jener} usually refers to something farther away or more distant in time. \textbf{jener} often sounds literary or formal.}
\end{grammarentry}

\begin{center}
\begin{tabular}{|>{\centering\arraybackslash}m{2.2cm}|>{\centering\arraybackslash}m{2.5cm}|>{\centering\arraybackslash}m{2.5cm}|>{\centering\arraybackslash}m{2.5cm}|>{\centering\arraybackslash}m{2.5cm}|}
\hline
\HeaderRow
\hline
\rowcolor{RowBlueLight}Nominativ & dieser / jener & diese / jene & dieses / jenes & diese / jene \\
\hline
\rowcolor{RowBlueDark}Akkusativ & diesen / jenen & diese / jene & dieses / jenes & diese / jene \\
\hline
\rowcolor{RowBlueLight}Dativ & diesem / jenem & dieser / jener & diesem / jenem & diesen / jenen \\
\hline
\rowcolor{RowBlueDark}Genitiv & dieses / jenes & dieser / jener & dieses / jenes & dieser / jener \\
\hline
\end{tabular}
\end{center}

\begin{exbox}
\begin{itemize}
\ExampleItem{Dieser Mann arbeitet hier.}{This man works here.}
\ExampleItem{An jenem Tag begann alles.}{Everything began on that day.}
\ExampleItem{In jenem Moment wusste ich es.}{At that moment, I knew it.}
\end{itemize}
\end{exbox}

\section{Der, die und das demonstrativ gebraucht}
\begin{grammarentry}{der / die / das}{demonstrativ gebraucht}
\GermanText{Die bestimmten Artikel \textbf{der}, \textbf{die} und \textbf{das} können auch demonstrativ gebraucht werden. Dann werden sie stärker betont und bedeuten ungefähr \glqq dieser / diese / dieses\grqq.}
\EnglishText{The definite articles \textbf{der}, \textbf{die}, and \textbf{das} can also be used demonstratively. They are then stressed more strongly and mean approximately ``this one'' or ``that one.''}
\end{grammarentry}

\begin{center}
\begin{tabular}{|>{\centering\arraybackslash}m{2.2cm}|>{\centering\arraybackslash}m{2.5cm}|>{\centering\arraybackslash}m{2.5cm}|>{\centering\arraybackslash}m{2.5cm}|>{\centering\arraybackslash}m{2.5cm}|}
\hline
\HeaderRow
\hline
\rowcolor{RowBlueLight}Nominativ & der & die & das & die \\
\hline
\rowcolor{RowBlueDark}Akkusativ & den & die & das & die \\
\hline
\rowcolor{RowBlueLight}Dativ & dem & der & dem & den \\
\hline
\rowcolor{RowBlueDark}Genitiv & dessen & deren & dessen & deren \\
\hline
\end{tabular}
\end{center}

\section{Derselbe}
\begin{grammarentry}{derselbe}{Identität: genau dieselbe Person oder Sache}
\GermanText{\textbf{derselbe}, \textbf{dieselbe} und \textbf{dasselbe} bedeuten, dass zwei Bezugnahmen genau dieselbe Person oder Sache meinen.}
\EnglishText{\textbf{derselbe}, \textbf{dieselbe}, and \textbf{dasselbe} mean that two references concern exactly the same person or thing.}
\end{grammarentry}

\begin{center}
\begin{tabular}{|>{\centering\arraybackslash}m{2.2cm}|>{\centering\arraybackslash}m{2.5cm}|>{\centering\arraybackslash}m{2.5cm}|>{\centering\arraybackslash}m{2.5cm}|>{\centering\arraybackslash}m{2.5cm}|}
\hline
\HeaderRow
\hline
\rowcolor{RowBlueLight}Nominativ & derselbe & dieselbe & dasselbe & dieselben \\
\hline
\rowcolor{RowBlueDark}Akkusativ & denselben & dieselbe & dasselbe & dieselben \\
\hline
\rowcolor{RowBlueLight}Dativ & demselben & derselben & demselben & denselben \\
\hline
\rowcolor{RowBlueDark}Genitiv & desselben & derselben & desselben & derselben \\
\hline
\end{tabular}
\end{center}

\section{Derjenige}
\begin{grammarentry}{derjenige}{Demonstrativpronomen, häufig mit Relativsatz}
\GermanText{\textbf{derjenige}, \textbf{diejenige}, \textbf{dasjenige} und \textbf{diejenigen} bezeichnen eine bestimmte Person oder Sache aus einer Gruppe. Danach folgt sehr häufig ein Relativsatz.}
\EnglishText{\textbf{derjenige}, \textbf{diejenige}, \textbf{dasjenige}, and \textbf{diejenigen} identify a particular person or thing from a group. They are very often followed by a relative clause.}
\end{grammarentry}

\begin{center}
\begin{tabular}{|>{\centering\arraybackslash}m{2.2cm}|>{\centering\arraybackslash}m{2.5cm}|>{\centering\arraybackslash}m{2.5cm}|>{\centering\arraybackslash}m{2.5cm}|>{\centering\arraybackslash}m{2.5cm}|}
\hline
\HeaderRow
\hline
\rowcolor{RowBlueLight}Nominativ & derjenige & diejenige & dasjenige & diejenigen \\
\hline
\rowcolor{RowBlueDark}Akkusativ & denjenigen & diejenige & dasjenige & diejenigen \\
\hline
\rowcolor{RowBlueLight}Dativ & demjenigen & derjenigen & demjenigen & denjenigen \\
\hline
\rowcolor{RowBlueDark}Genitiv & desjenigen & derjenigen & desjenigen & derjenigen \\
\hline
\end{tabular}
\end{center}

\begin{exbox}
\begin{itemize}
\ExampleItem{Diejenige, die dort steht, ist meine Schwester.}{The one who is standing there is my sister.}
\ExampleItem{Ich helfe demjenigen, der mich darum gebeten hat.}{I help the person who asked me to do so.}
\ExampleItem{Das sind diejenigen, die gewonnen haben.}{Those are the ones who won.}
\end{itemize}
\end{exbox}

\section{Solcher}
\begin{grammarentry}{solcher}{Hinweis auf eine Art oder Beschaffenheit}
\GermanText{\textbf{solcher}, \textbf{solche} und \textbf{solches} weisen auf eine bestimmte Art, Beschaffenheit oder Gruppe hin.}
\EnglishText{\textbf{solcher}, \textbf{solche}, and \textbf{solches} point to a certain type, quality, or group.}
\end{grammarentry}

\begin{center}
\begin{tabular}{|>{\centering\arraybackslash}m{2.2cm}|>{\centering\arraybackslash}m{2.5cm}|>{\centering\arraybackslash}m{2.5cm}|>{\centering\arraybackslash}m{2.5cm}|>{\centering\arraybackslash}m{2.5cm}|}
\hline
\HeaderRow
\hline
\rowcolor{RowBlueLight}Nominativ & solcher & solche & solches & solche \\
\hline
\rowcolor{RowBlueDark}Akkusativ & solchen & solche & solches & solche \\
\hline
\rowcolor{RowBlueLight}Dativ & solchem & solcher & solchem & solchen \\
\hline
\rowcolor{RowBlueDark}Genitiv & solchen & solcher & solchen & solcher \\
\hline
\end{tabular}
\end{center}

\section{Mancher}
\begin{grammarentry}{mancher}{Indefinitpronomen oder Artikelwort}
\GermanText{\textbf{mancher}, \textbf{manche} und \textbf{manches} bezeichnen eine unbestimmte einzelne Person oder Sache beziehungsweise mehrere unbestimmte Personen oder Sachen.}
\EnglishText{\textbf{mancher}, \textbf{manche}, and \textbf{manches} refer to an unspecified individual person or thing, or to several unspecified people or things.}
\GermanText{Grammatisch gehört \textbf{mancher} normalerweise zu den Indefinitpronomen und nicht zu den Demonstrativpronomen.}
\EnglishText{Grammatically, \textbf{mancher} normally belongs to the indefinite pronouns, not to the demonstrative pronouns.}
\end{grammarentry}

\begin{center}
\begin{tabular}{|>{\centering\arraybackslash}m{2.2cm}|>{\centering\arraybackslash}m{2.5cm}|>{\centering\arraybackslash}m{2.5cm}|>{\centering\arraybackslash}m{2.5cm}|>{\centering\arraybackslash}m{2.5cm}|}
\hline
\HeaderRow
\hline
\rowcolor{RowBlueLight}Nominativ & mancher & manche & manches & manche \\
\hline
\rowcolor{RowBlueDark}Akkusativ & manchen & manche & manches & manche \\
\hline
\rowcolor{RowBlueLight}Dativ & manchem & mancher & manchem & manchen \\
\hline
\rowcolor{RowBlueDark}Genitiv & manchen & mancher & manchen & mancher \\
\hline
\end{tabular}
\end{center}

\section{Zusammenfassung}
\begin{grammarentry}{Merksatz}{kurz und wichtig}
\GermanText{\textbf{dieser / jener}, \textbf{derselbe} und \textbf{derjenige} können demonstrativ gebraucht werden. \textbf{mancher} ist dagegen normalerweise ein Indefinitpronomen.}
\EnglishText{\textbf{dieser / jener}, \textbf{derselbe}, and \textbf{derjenige} can be used demonstratively. \textbf{mancher}, by contrast, is normally an indefinite pronoun.}
\GermanText{Alle diese Formen werden nach Kasus, Genus und Numerus dekliniert.}
\EnglishText{All these forms are declined according to case, gender, and number.}
\end{grammarentry}

\end{document}
